%% kapa: exam=1 round=21 year=2010 period=1 form=A subject=civil_law subject_ko=민법 seq=37 number=37
%% source: Q-Net 감정평가사 기출문제 (cst00309) artlSeq=5024181 file=extracted/1008/5024181/2010년 제21회 감정평가사 1차시험 기출문제/1교시.hwp
%% conversion: kordoc 4.24.51 -> tools/gen1.py (auto); review=proofread-sonnet; flags=none
\begin{question}{37}
하나의 부동산에 설정된 저당권과 용익물권의 관계에 관한 설명으로 옳지 않은 것은? (다툼이 있으면 판례에 의함)
\begin{choices}
  \item 1번 저당권이 설정된 후 지상권이 설정되고 그 후 2번 저당권이 설정된 경우, 2번 저당권 실행으로 목적물이 매각되더라도 지상권은 소멸하지 않는다.
  \item 전세권이 저당권보다 먼저 설정된 경우, 저당권 실행시 전세권자가 기한의 이익을 포기하고 배당요구를 하면 전세권은 목적물의 매각으로 소멸한다.
  \item 지상권이 저당권보다 먼저 설정된 경우, 저당권 실행으로 토지가 매각되더라도 지상권은 소멸하지 않는다.
  \item 전세권이 저당권보다 후에 설정된 경우, 전세권자가 목적물에 유익비를 지출하였다면 전세권자는 저당목적물의 매각대금에서 그 비용을 우선상환 받을 수 있다.
  \item 지상권이 저당권보다 후에 설정된 경우, 지상권자는 저당권자에게 그 토지로 담보된 채권을 변제하고 저당권의 소멸을 청구할 수 있다.
\end{choices}
\end{question}
\answer{1}
